\documentclass[11pt]{article}
\usepackage[utf8]{inputenc}
\usepackage{amsmath,amssymb,amsthm}
\usepackage[margin=1in]{geometry}
\usepackage{microtype}
\usepackage{hyperref}
\usepackage{xcolor}
\usepackage{enumitem}
\usepackage{newunicodechar}
\emergencystretch=3em
\newunicodechar{ℝ}{\ensuremath{\mathbb{R}}}
\newunicodechar{λ}{\ensuremath{\lambda}}
\newunicodechar{·}{\ensuremath{\cdot}}
\newunicodechar{∂}{\ensuremath{\partial}}

\newcommand{\tideref}[2][]{\href{https://therisensea.org/tides/#2/}{\ifx&#1&\texttt{#2}\else#1\fi}}

\title{Tide retrospective:\\\emph{linear response at any perturbation strength}}
\author{Claude Opus 4.8 (1M ctx), with Daniel Murfet}
\date{2026-05-30}

\begin{document}
\maketitle

\begin{abstract}
The susceptibility of the anharmonic Gibbs measure --- the derivative of the
perturbed mean $\langle u\rangle_h = G_1(h)/G_0(h)$ of $u=-(tx)$ --- lifted from
the single point $h=0$ to an arbitrary base point $|h_0|<1$:
\[
  \frac{\partial}{\partial h}\Big|_{h_0}\langle u\rangle_h
    = \frac{G_2(h_0)G_0(h_0)-G_1(h_0)^2}{G_0(h_0)^2}.
\]
This is the fluctuation--dissipation (linear-response) identity at any
perturbation strength, and the foundation for the third cumulant (the second
derivative of the mean). New file
\texttt{AnharmonicSusceptibilityGeneralH.lean}; builds clean first try, no
\texttt{sorry}/\texttt{axiom}/\texttt{native\_decide}.
\end{abstract}

\section{Setting}

Seventh tide of a single-day anharmonic arc on the laplace seabed. The
previous tide\footnote{\tideref{2026-05-30-tide-anharmonic-second-cumulant}.}
proved the susceptibility $\partial_h\langle u\rangle_h$ only at $h=0$ (the
connected variance $\kappa_2$). The third cumulant is the \emph{second}
derivative of the mean, so it needs the mean's derivative as a function of
$h$ on a neighbourhood --- which in turn needs the susceptibility at a general
base point. This tide provides exactly that, at minimal risk: the general-$h$
partition derivatives already exist (from the general-$h$ step tide), so the
only genuinely new pieces are positivity of the perturbed partition at $h_0\neq0$
and a pointwise quotient rule.

\section{The thing formalised}

\begin{quote}\itshape
For the anharmonic potential and $|h_0|<1$, $\mathrm{HasDerivAt}$ of
$h\mapsto G_1(h)/G_0(h)$ at $h_0$ with value
$(G_2(h_0)G_0(h_0)-G_1(h_0)^2)/G_0(h_0)^2$; and $G_0(h)>0$ for $|h|<1$.
\end{quote}

\section{Proof strategy}

\textbf{Positivity.} $G_0(h)=\int e^{-t(L+hx)}$ (after $(-(tx))^0=1$) is a
strictly-positive integrand with full support and is integrable at any
$|h|<1$ (the arc's centre-integrability lemma at $n=0$); so
\texttt{integral\_pos\_iff\_support\_of\_nonneg\_ae} gives positivity --- a
line-for-line analogue of \texttt{anharmonic\_partition\_pos} with the $+hx$
carried through. \textbf{Susceptibility.} \texttt{HasDerivAt.div} on the two
general-$h$ partition derivatives --- of $G_1$ (value $G_2(h_0)$) and of $G_0$
(value $G_1(h_0)$) at $h_0$ --- with $G_0(h_0)\neq0$.

\section{Roadblocks and resolutions}

None --- it built on the first attempt. The general-$h$ machinery from the
earlier tides (centre-integrability, the general-$h$ \texttt{HasDerivAt} step)
is exactly the right interface; positivity reused the $h=0$ template verbatim
modulo the perturbation term.

\section{What was Mathlib, what was new}

\textbf{Mathlib.} \texttt{HasDerivAt.div}; the positivity-from-support
lemma\footnote{\texttt{integral\_pos\_iff\_support\_of\_nonneg\_ae}.};
\texttt{Real.exp\_ne\_zero}; \texttt{Real.volume\_univ}.

\textbf{Seabed (reused).} The general-$h$ partition derivative and the
centre-integrability lemma\footnote{\texttt{weightedPartition\_hasDerivAt}
and \texttt{integrable\_weightedPartition\_integrand}.}, plus the $h=0$
positivity template --- all from this session's arc.

\textbf{New.} General-$h$ partition positivity and the general-$h$
susceptibility. About 70 lines.

\section{Lessons}

\begin{itemize}[itemsep=3pt]
\item \textbf{Investing in the general-$h$ step pays compounding dividends.}
  Because the partition derivatives were proved at an arbitrary centre two
  tides ago, the general-$h$ susceptibility is just ``the $h=0$ proof with
  $h_0$ in place of $0$'' --- no new analytic work, only the (reused)
  positivity and a quotient rule.
\item \textbf{Positivity templates transfer across perturbations.} The
  $\int e^{-tL}>0$ argument carries the $+hx$ term unchanged: the integrand is
  still everywhere positive with full support.
\end{itemize}

\section{Follow-ups}

\begin{itemize}[itemsep=3pt]
\item \textbf{Third cumulant $\kappa_3$.} The second derivative of the mean:
  differentiate the general-$h$ susceptibility's mean function once more (the
  \texttt{iteratedDeriv} chaining pattern), then expand to
  $\langle u^3\rangle_0-3\langle u^2\rangle_0\langle u\rangle_0+2\langle u\rangle_0^3$.
  Then $\kappa_4$ --- the flow-equation identity (survey v4 \#2).
\item \textbf{Promote \texttt{amgm\_t\_abs\_x} to \texttt{lean-common}.}
\end{itemize}

\end{document}
